\section{Pairings and the positive-root restriction}

An ambient Hopf pairing and an invariant form on a primitive Lie algebra
have different compatibility equations.  Their distinction becomes
visible on a product of two primitives.

Let $C,D$ be ordinary connected positively graded bialgebras.
An indexed pairing is a family
$\langle-,-\rangle_n:C_n\otimes D_n\to\kk$, where the second index
is interpreted as dual height.  Extend it to $C\times D$ by zero on
unequal heights.  Suppose it is normalized by
\[
 \langle1,z\rangle=\varepsilon_D(z),\qquad
 \langle x,1\rangle=\varepsilon_C(x),
\]
and satisfies the ambient product--coproduct identity
\begin{equation}
 \langle xy,z\rangle
 =\langle x\otimes y,\Delta_Dz\rangle,
 \quad
 \langle x\otimes y,u\otimes v\rangle
 =\langle x,u\rangle\langle y,v\rangle.
 \label{ambient-adjointness}
\end{equation}
For a super pairing the tensor pairing includes the Koszul interchange.
A pairing defined only on primitives does not type
\eqref{ambient-adjointness}, since neither $xy$ nor the two factors of
$\Delta z$ need be primitive.

\begin{proposition}[The primitive restriction of a Hopf pairing]
\label{primitive-pairing}
Under these hypotheses,
\[
 \langle[P(C),P(C)],P(D)\rangle=0.
\]
Thus the quotient of $P(C)$ by the left radical of its pairing with
$P(D)$ is abelian.  The same conclusion holds with supercommutators
and a normalized super pairing.
\end{proposition}

\begin{proof}
For primitives $x,y\in C$ and $z\in D$, the right side of
\eqref{ambient-adjointness} is
\[
 \langle x,z\rangle\varepsilon_C(y)
 +\varepsilon_C(x)\langle y,z\rangle=0.
\]
Reversing $x,y$ gives zero again, so their commutator pairs to zero.
In the super case the possible interchange signs multiply zero terms.
The left radical contains the entire derived Lie algebra, so it is a
Lie ideal and its quotient is abelian.
\end{proof}

Consequently, a nonabelian Borcherds recognition cannot obtain its
nondegenerate primitive form by restricting this same ambient pairing.
An invariant Lie form needs its own construction.  A Hall double may
use pairings between separately specified positive and negative halves,
with a proved comparison to the Lie form.

There are two further type restrictions.  A morphism
$C\bar\otimes_ND\to\kk$ in the positive-height category vanishes on
every positive total height.  It cannot realize a nonzero indexed
root pairing.  Also, indexed pairings define the evaluation
\[
 \prod_{n\geq0}C_n\ \times\ \bigoplus_{n\geq0}D_n
 \longrightarrow\kk
\]
by a finite sum.  They need not define a pairing between two products.
If a bilinear extension between the two products were jointly continuous
over discrete $\kk$, it would vanish on $F^r\widehat C\times
F^s\widehat D$ for some $r,s$.
It would then vanish on all paired components of height
$n\geq\max(r,s)$.  This contradicts a nonzero diagonal pairing in
infinitely many heights.  A continuous/discrete dual pairing or an
explicit summability condition is therefore part of a completed double.

\begin{proposition}[Even primitives in ordinary finite bialgebras]
\label{finite-primitive-obstruction}
A finite-dimensional super bialgebra over a characteristic-zero field
has no nonzero even primitive.  Its primitive Lie superalgebra is abelian.
\end{proposition}

\begin{proof}
Let $x$ be such a primitive, and let $f$ be its monic minimal polynomial,
of degree $n$.  A primitive scalar is zero, since applying the counit to
its primitive equation gives zero counit.  Hence $n\geq2$.
The vectors $1,x,\ldots,x^{n-1}$ are linearly independent.
As $x$ is even,
\[
 0=\Delta f(x)=f(x\otimes1+1\otimes x).
\]
In this tensor basis the coefficient of $x^{n-1}\otimes x$ is $n$.
Reducing the two pure $x^n$ terms contributes only tensors having a unit
factor.  Lower terms of $f$ have total exponent at most $n-1$ and cannot
contribute.  Thus $n=0$ in the field, a contradiction.
Every odd--odd bracket is an even primitive, so these brackets vanish
as well.  With no even primitives, this proves the second assertion.
\end{proof}

This obstruction applies to ordinary tensor-valued finite bialgebras.
The total-height windows of Remark~\ref{polynomial-window} retain an
even primitive because their bialgebra equation is imposed in a different
tensor category.  In characteristic $p$, the ordinary algebra
$\kk[x]/(x^p)$ does admit the primitive coproduct: all intermediate
binomial coefficients in $\Delta(x^p)$ vanish.
An odd exterior generator also gives a finite-dimensional super Hopf
algebra in characteristic zero.
The monoidal restriction, tower reconstruction, antipode, and
differential-stable quotient proofs themselves hold over every field.
The even-primitive obstruction and the enveloping primitive comparison
use characteristic zero.

The positive cone itself is necessary for this truncation construction.
More generally, it works for a commutative monoid $M$ with an additive
height $h:M\to\mathbb N$, satisfying $h^{-1}(0)=\{0\}$ and finite
$h^{-1}([0,N])$.  These assumptions make retained sets closed under
summands and make each charge decomposition finite.
An arbitrary finite subset of a charge group has no such property.
Retaining zero and all summands of a retained sum would require retaining
every $\alpha$, since $\alpha+(-\alpha)=0$.

For a real-root $\mathfrak{sl}_2$ triple, the relations
\[
 [h_i,e_i]=2e_i,\qquad[e_i,f_i]=h_i
\]
force $h(h_i)=0$ and then $h(e_i)=h(f_i)=0$ under any nonnegative
additive height making these nonzero elements homogeneous.
They cannot be nonzero primitives in a connected height-zero component
$\kk1$.  Allowing a nontrivial degree-zero Hopf algebra removes this
connectedness obstruction.  It does not give finite-dimensional windows
containing an even Cartan primitive, by
Proposition~\ref{finite-primitive-obstruction}.
Thus the finite-height recognition theorem concerns the positive Lie
algebra.  A full root algebra requires separately specified halves,
Cartan data, and a compatible double or another completion.
